\documentclass[hyperref={colorlinks=true,linkcolor=blue}, aspectratio=169]{beamer}

\usepackage{animate}

\usepackage{booktabs}
\setbeamertemplate{caption}[numbered]
\setbeamertemplate{footline}[frame number]
\setbeamertemplate{navigation symbols}{}

\usepackage{amsmath}
\newcommand*{\norm}[1]{\left\lVert#1\right\rVert}
\newcommand{\argmin}{\arg\!\min}

\usepackage{color,soul}
\usepackage{colortbl}

\setbeamerfont{footnote}{size=\tiny}

\usepackage{listings,lstautogobble}
\lstset{language=Python,
	basicstyle=\scriptsize\ttfamily,
	commentstyle=\ttfamily\itshape\color{gray},
	showstringspaces=false,
	autogobble=true
}

\usepackage{caption}
\captionsetup{justification=centering,margin=1cm}
\usepackage{subcaption}

\usepackage{multirow}

\usepackage{siunitx}
\usepackage{tikz}
\tikzstyle{process} = [rectangle, minimum width=3cm, minimum height=1cm, text centered, draw=black, fill=orange!30]
\usetikzlibrary{shapes,positioning}
\usepackage{ulem}

\AtBeginSection[]{
	\begin{frame}
		\vfill
		\centering
		\begin{beamercolorbox}[sep=8pt,center,shadow=true,rounded=true]{title}
			\usebeamerfont{title}\insertsectionhead\par%
		\end{beamercolorbox}
		\vfill
	\end{frame}
}

\addtobeamertemplate{footnote}{\hskip -2em}{}

\usepackage[capitalise,noabbrev]{cleveref}
\usepackage[numbers,sort&compress]{natbib}
\bibliographystyle{unsrt}

\setbeamercolor{block title}{fg=black,bg=yellow}

\title{On the use of autoencoder to denoise diffusion MRI}
\author[Z Tan]{Soundarya Soundarresan \inst{1}, \underline{Zhengguo Tan} \inst{1}, Patrick A Liebig \inst{2}, \\ Robin M Heidemann \inst{2}, Frederik B Laun \inst{3}, Florian Knoll \inst{1}}
\institute{\inst{1} Artificial Intelligence in Biomedical Engineering, University of Erlangen-Nuremberg, Erlangen, Germany \\
	\inst{2} Siemens AG, Erlangen, Germany \\
	\inst{3} Institute of Radiology, University Hospital Erlangen, Erlangen, Germany}
\date{05 October 2023}

\begin{document}
	
	\frame{\titlepage}
	
	% ================= %
	\section{Motivation: In vivo High Resolution Diffusion MRI at 7~T}
	
	\begin{frame}{Various clinical applications of diffusion MRI 
			\footnote{Courtesy: ChatGPT}$^,$
			\footnote{\url{https://mriquestions.com/trace-vs-adc-map.html}}$^,$ \footnote{\url{https://www.psychologytoday.com/gb/blog/adventures-in-cognition/202102/why-white-matter-matters-diffusion-tensor-imaging}}$^,$
			\footnote{\url{https://rad.bwh.harvard.edu/prostate-mri/}}}
		\begin{columns}
			\begin{column}{0.5\textwidth}
				Brain: stroke \& neurodegenerative diseases 
				
				\vspace{1em}
				
				\includegraphics[height=0.28\textheight]{figures/diffusion.png}
				
				
			\end{column}
			\vline
			\hspace{0.1em}
			\begin{column}{0.5\textwidth}
				Prostate cancer: detection and localization 
				
				\vspace{1em}
				
				\includegraphics[height=0.28\textheight]{figures/PI_RAD.png}
			\end{column}
		\end{columns}

	\end{frame}

	\begin{frame}{Limitations in diffusion MRI \footnote{Le Bihan D, et al. Artifacts and pitfalls in diffusion MRI. \textit{J Magn Reson Imaging} (2006).}}
		\begin{columns}
			\begin{column}{0.6\textwidth}
				\begin{itemize}
					\item Common pitfalls:
					\begin{itemize}
						\item [$\star$] motion artifacts
						\item [$\star$] eddy currents
						\item [$\star$] geometric distortion and susceptibility artifacts
					\end{itemize}
					\vspace{2em}
					\item Be aware that diffusion-weighted imaging (DWI),
					\begin{itemize}
						\item [$\star$] primarily uses single-shot echo planar imaging (EPI)
						\item [$\star$] clinical spatial resolution of about \SI{2}{mm}
						\item [$\star$] suffers from low signal-to-noise ratio (SNR)
					\end{itemize}
				\end{itemize}
			\end{column}

			\begin{column}{0.4\textwidth}
				\includegraphics[height=0.8\textheight]{figures/diffusion_pitfalls.png}
			\end{column}
		\end{columns}

	\end{frame}

	% ================= %
	\section{Method}

	\begin{frame}{Diffusion MRI Acquisition}
		\begin{figure}
			\centering
			\includegraphics[width=0.7\textwidth]{figures/diffusion_sampling.png}
		\end{figure}

		\begin{itemize}
			\item <1-> 2-shot interleaved EPI with multi-band factor 2, yielding acceleration factor $6 \times 2$
			\item <2-> $k_y$-shift encoding to accomplish complementary $k$-$q$-space sampling
			\item <3-> \SI{1.2}{mm} (isotropic), $2 \times b_0$ + 30 diffusion encodings with $b$-value \SI{1000}{\second\per\square\mm}
		\end{itemize}
	\end{frame}

	\begin{frame}{Diffusion MRI Forward Modeling and Image Reconstruction \footnote{Tan Z, et al. \textit{under review}}}
		\begin{figure}
			\centering
			\includegraphics[width=\textwidth]{figures/fwd.pdf}
		\end{figure}
	\end{frame}

	\begin{frame}{Denoise Diffusion MRI Using Autoencoder \footnote{Hinton GE, Salakhutdinov RR. Reducing the dimensionality of data with neural networks. \textit{Science} (2006).}$^,$
		\footnote{Lam F, et al. Constrained magnetic resonance spectroscopic imaging by learning nonlinear low-dimensional models. \textit{IEEE Trans Med Imaging} (2019).}$^,$
		\footnote{Mani M, et al. qModeL: A plug-and-play model-based reconstruction for highly accelerated multi-shot diffusion MRI using learned priors. \textit{Magn Reson Med} (2021).}}
		\begin{figure}
			\centering
			\includegraphics[width=0.8\textwidth]{figures/dae.pdf}
		\end{figure}
	\end{frame}

	\begin{frame}{Denoise Diffusion MRI Using Autoencoder: How to Train the Model}
		\begin{itemize}
			\item Data generation:

			\begin{enumerate}
				\item Simulation of DWI signal based on the DTI model \footnote{Basser PJ, et al. MR diffusion tensor spectroscopy and imaging. \textit{Biophys J} (1994).}
				\item Each tensor element contains 10 entries, resulting in $10^6$ DWI signal entries
				\item Zero-mean Gaussian-distribution noise with 10 varying standard deviations are added to the simulated signal, yielding $10^7$ DWI signal
				\item \SI{80}{\percent} of the data for training, and \SI{20}{\percent} for testing
			\end{enumerate}
			\vspace{1em}
			\item Training setup:
			\begin{enumerate}
				\item Optimizer: ADAM \footnote{Kingma DP, et al. Adam: A Method For Stochastic Optimization. \textit{ICLR} (2015).}
				\item Learning rate: 0.0001
				\item Epochs: 800
				\item 4 fully-connected encoder/decoder layers with ReLU, whereas the output layer uses Sigmoid as the activation function
				\item Latent size: 10
			\end{enumerate}
		\end{itemize}
	\end{frame}
	
	% ================= %
	\section{Results}

	\begin{frame}{Comparison with the Linear Subspace Method}
		\begin{figure}
			\includegraphics[height=0.8\textheight]{figures/SVD_DAE.png}
		\end{figure}
	\end{frame}

	\begin{frame}{In vivo DWI Denoising}
		\begin{figure}
			\includegraphics[width=0.7\textwidth]{figures/MUSE_train_DAE_Latent10_Depth04_epoch0100_dwi.png}
		\end{figure}
	\end{frame}

	\begin{frame}{In vivo DWI Denoising: Learned Latent Signal}
		\begin{figure}
			\includegraphics[width=\textwidth]{figures/MUSE_train_DAE_Latent10_Depth04_epoch0100_latent.png}
		\end{figure}
	\end{frame}

	% ================= %
	\section{Conclusion}
	
	\begin{frame}{Conclusion}
		\begin{columns}
			\begin{column}{0.5\textwidth}
				\begin{itemize}
					\item Multi-shot EPI:
					\begin{itemize}
						\item [$\star$] enables high spatial resolution DWI
					\end{itemize}
					\vspace{1em}
					\item we open source our DWI reconstruction techniques $\rightarrow$
					\vspace{1em}
					\item Autoencoder:
					\begin{itemize}
						\item [$\star$] enables DWI denoising
						\item [$\star$] learns a nonlinear subspace 
					\end{itemize}
				\end{itemize}
			\end{column}
			
			\begin{column}{0.5\textwidth}
				\includegraphics[width=\columnwidth]{figures/github_jets.png}
			\end{column}
		\end{columns}
		
	\end{frame}
	
\end{document}